\begin{document}

{
\setcounter{tocdepth}{2}
\tableofcontents
}
\section{TFPT / Universalraum: die richtige einfache
Kette}\label{tfpt-universalraum-die-richtige-einfache-kette}

\subsection{Vollständige Konsolidierung v1.6.7 - 15. September
2026}\label{vollstuxe4ndige-konsolidierung-v1.6.7---15.-september-2026}

\subsubsection{Ergebnis und
Leseschlüssel}\label{ergebnis-und-leseschluxfcssel}

Die zuletzt ergänzte Vogelperspektive steht im Teil B dieser Fassung:
\textbf{Gesucht wird der kleinste gemeinsame beobachtbare Prozess hinter
den verschiedenen Ansichten.} Ein nativer Rang-60-Kodierer ist bereits
vorhanden. Seine Informationsverluste, die minimalen
symmetrieverträglichen Dynamiken und der Unterschied zwischen
Korrelation und kausalem Eingriff wurden neu berechnet. Das verschiebt
die fundamentale Priorität von weiteren lokalen Näherungswerten zur
Herkunft der Operationen und zur gemeinsamen Schattenkarte.

Die neue Runde schließt zwei kleine Strukturfragen vollständig: Der
symmetrische Zustandsraum bei zwei Bosonen hat exakt vier Dimensionen;
der bisher unbestimmte Rest von 4035 bilinearen Komponenten zerfällt
vollständig in drei Darstellungen. Die vier Zustandsrichtungen und ihre
Kopplungen wurden zusätzlich am ursprünglichen Tensor berechnet. Daraus
folgt ein vollständiger sechs-dimensionaler Ausschnitt.

Die gewünschte einfache Kette existiert ebenfalls, aber anders als
vorgeschlagen: Man muss den \textbf{ganzen Hamiltonoperator}
wiederholen, nicht Bosonzahl und Iterationsnummer gleichsetzen. Die
ersten fünf Lanczos-Glieder sind jetzt exakt bekannt. Sie bestimmen zehn
Energiemomente des gefüllten Startzustands und verschärfen die obere
Grundzustandsgrenze geringfügig.

\textbf{Das ist keine vollständige TOE.} Insbesondere fehlen die
physische Auswahl des Modellvertrags, ausführbare ursprüngliche
Instrumente, derselbe geladene Transport im gemeinsamen Träger und die
relativistische Kontinuumskonstruktion. Dass die bisherigen Grunddaten
dafür noch nicht eindeutig genügen, wird unten an einer einfachen
Familie verschiedener Dynamiken nachgewiesen.

Kennzeichnungen: \textbf{exakt} bezeichnet eine bewiesene Aussage im
genannten Modell; \textbf{bedingt} eine Aussage mit zusätzlich gewährter
Operation oder bekanntem früherem Satz; \textbf{numerisch} eine
Näherung; \textbf{offen} einen fehlenden Nachweis. Komplexe
Algebradimensionen zählen weder physikalische Raumrichtungen noch
unabhängig hergeleitete Teilchen oder ausführbare Operationen.

Die Fassung baut auf dem vollständigen Hauptdokument v1.6.6 auf. Die
aktuelle Gesamtdarstellung steht voran; die frühere Fassung
einschließlich ihrer Herleitungen, Gegenbeispiele und älteren Anhänge
bleibt vollständig erhalten. Alle sechs in dieser Runde eingesandten
Texte werden unverändert dokumentiert und hier mit überprüften
Ergänzungen und ausdrücklich benannten Korrekturen gelesen.

\subsection{1. Das Gesamtbild: Grammatik, Prozess, Zustand,
Antwort}\label{das-gesamtbild-grammatik-prozess-zustand-antwort}

TFPTs Compiler beschreibt zunächst eine markierte algebraische
Grammatik: welche Bausteine zusammenpassen, welche Symmetrien gelten und
welche Zahlen innerhalb dieser Konstruktion folgen. Der Universalraum
soll diese Grammatik als einen ausführbaren Prozess realisieren. Der
entscheidende Unterschied ist der zwischen einem erlaubten Zusammenhang
und einer tatsächlich erzeugten physikalischen Entwicklung.

Im hier untersuchten Modell ist die Grundregel weiterhin sehr einfach:
\textbf{Zwei Fermionen können in einen Vermittler umgewandelt werden und
zurück.} Mit 64 Fermionmoden und 60 Bosonmoden lautet sie

\[\adjustbox{max width=\linewidth}{$\displaystyle 
H=\Delta N_b+g(T_++T_-),\qquad
T_+=\sum_A b_A^\dagger P_A,\quad T_-=T_+^\dagger,
\qquad P_A=\sum_{i<j}W_{A,ij}f_jf_i.
$}\]

Es gelten \(\Delta>0\), reelles \(g\) und die erhaltene Ladung
\(N=N_f+2N_b\). Der Tensor W besitzt 480 Einträge mit Werten +1 oder -1,
acht disjunkte Paare pro Zeile, und \(WW^\dagger=8I_{60}\). Seine innere
Darstellung ist \((16,4)\) für Fermionen und \((10,6)\) für Bosonen
unter Spin(10) mal SU(4).

Die heutige Beweiskette sieht so aus:

{\def\LTcaptype{none} % do not increment counter
\begin{longtable}[]{@{}
  >{\raggedright\arraybackslash}p{(\linewidth - 4\tabcolsep) * \real{0.3333}}
  >{\raggedright\arraybackslash}p{(\linewidth - 4\tabcolsep) * \real{0.3333}}
  >{\raggedright\arraybackslash}p{(\linewidth - 4\tabcolsep) * \real{0.3333}}@{}}
\toprule\noalign{}
\begin{minipage}[b]{\linewidth}\raggedright
Verbindung
\end{minipage} & \begin{minipage}[b]{\linewidth}\raggedright
Was bereits trägt
\end{minipage} & \begin{minipage}[b]{\linewidth}\raggedright
Was dadurch nicht automatisch folgt
\end{minipage} \\
\midrule\noalign{}
\endhead
\bottomrule\noalign{}
\endlastfoot
Compiler zu Paarregel & Konkreter Tensor, Symmetrie und Ladungserhaltung
& Physische Wahl von Energie, Kopplung und Instrumenten \\
Paarregel zu Zustand & Früherer eindeutiger Grundzustandssatz im
festgelegten Vertrag & Herleitung dieses Vertrags oder dessen
Präparation \\
Zustand zu geladener Antwort & Isolierte niedrige Fermion-Entnahmelinie
mit hohem Gewicht & Räumlich laufendes relativistisches Teilchen \\
Antwort zu Bewegung & Zusätzliche Links erlauben nachgewiesenen Transfer
& Dass derselbe Link aus dem ursprünglichen Operationssatz entsteht \\
Bewegung zu Raumzeit & Prüfprogramm und bedingte Modelle & Gemeinsame
3+1D-Welt, chirales Maß und dynamische Gravitation \\
\end{longtable}
}

Der endliche dokumentierte Clock gehört zur inneren Spin(10)-Symmetrie.
Er ist damit kein zusätzliches äußeres Symmetrieobjekt. Seine Periode
sechs ist aber auch noch keine Herleitung einer physikalischen
Zeiteinheit oder einer Raumzeit. Innere Spin(10)-Labels sind nicht von
selbst Lorentzspin.

\subsection{2. Was die beiden neuen Texte
beitragen}\label{was-die-beiden-neuen-texte-beitragen}

\subsubsection{2.1 Quellenstand und
Reproduktionsgrenzen}\label{quellenstand-und-reproduktionsgrenzen}

17 ausgewählte Eingaben wurden mit SHA-256 eingefroren: beide
Nutzertexte, die zugehörigen Programme und Ergebnisdateien, der native
Tensor sowie die vollständige frühere Fassung mit Prüfpaket. Keine
fremde Quelldatei wurde verändert. Die Herkunftsdateien gehören zu
parallel fortgeschriebenen Arbeitsständen; Dateipins sind deshalb
wichtiger als Versionsnamen allein.

Der neue Grundzustandslauf war beim ersten Blick noch ohne
Ergebnisdatei. Beim Einfrieren lag bereits sein fertiges Ergebnis vor:
\textbf{382 Guards, davon 379 exakt und drei numerisch}. Sein
Quellenhash stimmt mit der eingefrorenen Programmdatei überein. Er
findet ebenfalls die Multiplizitäten 1, 1 und 4. Der Zwischenhinweis
dieser Prüfung auf eine noch leere Datei war bei der späteren Übernahme
bereits überholt und wird hier korrigiert.

Das Operations-JSON enthält \textbf{306}, nicht die im Begleittext
genannten 307 Guards; sein interner Checkerhash passt nicht zur
gleichzeitig eingefrorenen Programmdatei. Das ist ein Versionskonflikt
des Berichts, kein Beweis, dass seine Mathematik falsch ist. Die
tragenden Aussagen werden deshalb unabhängig geprüft, nicht als exakter
Gesamtreplay dieses fremden Programms ausgegeben. Das Feld-JSON meldet
2161 Guards und besitzt einen passenden Quellenhash; es enthält dennoch
den unten korrigierten Normfehler und zu weit formulierte
Interpretationen.

Mit den weiteren Texten und zugehörigen Auditdateien sind es insgesamt
24 eingefrorene Eingaben. Unsere eigene neue Suite hat \textbf{822
exakte Prüfbedingungen pro Python-Modus}. Normaler und optimierter Lauf
erzeugen identische Ergebnisdateien; Warnungen werden als Fehler
behandelt. Zusätzlich wurde die ältere ganzzahlige Kontraktionsrechnung
frisch kompiliert und wiederholt, mit identischen Zahlen. Die
vier-Boson-Norm aus dem früheren Zertifikat bleibt ein gekennzeichneter
Eingang: Ihre vollständige Enumeration wurde in dieser Runde nicht
wiederholt. Auch die gesamten fremden Programme mit ihren Guards wurden
nicht alle erneut ausgeführt. Prüfzahl und Beweisumfang werden
ausdrücklich getrennt.

\subsubsection{2.2
Übernahmeentscheidung}\label{uxfcbernahmeentscheidung}

{\def\LTcaptype{none} % do not increment counter
\begin{longtable}[]{@{}
  >{\raggedright\arraybackslash}p{(\linewidth - 2\tabcolsep) * \real{0.5000}}
  >{\raggedright\arraybackslash}p{(\linewidth - 2\tabcolsep) * \real{0.5000}}@{}}
\toprule\noalign{}
\begin{minipage}[b]{\linewidth}\raggedright
Aussage aus den Eingaben
\end{minipage} & \begin{minipage}[b]{\linewidth}\raggedright
Entscheidung für v1.6.7
\end{minipage} \\
\midrule\noalign{}
\endhead
\bottomrule\noalign{}
\endlastfoot
Korrigierter Anfangskommutant 3829536 / 1444233216 & Bestätigt; keine
neue Größe gegenüber v1.6.6 \\
Nur SU(4) oder nur Spin(10) als zusätzliche Kontrollen & Übernommen;
echte Zwischenstufen statt binärer Verfügbarkeit \\
Volle Gruppe plus alle Modenbesetzungen erzeugt volle Sektoralgebra &
Unabhängig bestätigt, ausdrücklich bedingt \\
Clock außerhalb der zusammenhängenden Gruppe, weil nicht skalar &
Verworfen; das explizite Spin-Wort aus v1.6.6 widerlegt den Schluss \\
Vier Singuletts auf der zweiten Bosonstufe & Unabhängig bestätigt und
durch konkrete Zustandsprojektionen ergänzt \\
Eine Zustandsrichtung pro Bosonzahl könnte bis Stufe 32 genügen &
Verworfen; schon Stufe zwei hat vier Richtungen \\
Feldnorm 24 bei quadrierter Norm 192 & Korrigiert: Norm ist 8 mal Wurzel
3 in derselben Konvention \\
4035 bilineare Komponenten unbestimmt & Jetzt vollständig in drei innere
Darstellungen zerlegt \\
Jeder relativistische Abschluss müsse den Tensorfeldkanal verwenden &
Nur unter dem festgelegten gleichhändigen, ableitungsfreien Ansatz
gültig \\
Schwächere Sektorschranken trennen N=64 nicht & Richtig für diese Sonde;
widerlegt den stärkeren früheren Satz nicht \\
\end{longtable}
}

Die Entropie von ungefähr 1,66 Bit im ersten Text gehört zum dortigen
Ritz-Näherungszustand. Die Blockdiagonalität der reduzierten Bosondichte
rechtfertigt die Shannon-Untergrenze für diesen Zustand, nicht ungeprüft
denselben Zahlenwert für den wahren nativen Grundzustand. Ebenso ist
eine Kreuzung mit dem leeren Zustand bei einem bestimmten chemischen
Potential kein Beweis für die vollständige Sektor-Stabilitätszone.

\subsection{3. Vollständig gelöst: der erste verzweigte
Singulettbereich}\label{vollstuxe4ndig-geluxf6st-der-erste-verzweigte-singulettbereich}

\subsubsection{3.1 Warum vier und nicht
eins?}\label{warum-vier-und-nicht-eins}

Im Ladungssektor N=64 schreiben wir Zustände relativ zur gefüllten
Fermionreferenz \(F\). Bei k Bosonen gibt es 2k Löcher. Der
Singulettbereich auf dieser Stufe liegt in

\[\adjustbox{max width=\linewidth}{$\displaystyle 
\mathcal S_k=
\left[\Lambda^{2k}(\overline{16\otimes4})
\otimes\operatorname{Sym}^k(10\otimes6)\right]^{G},
\qquad G=\operatorname{Spin}(10)\times SU(4).
$}\]

Die ersten Dimensionen sind \textbf{1, 1, 4}. Dies ist eine Aussage über
alle Singuletts dieser Stufen, nicht über vier verschiedene
Grundzustände. Ein eindeutiger Grundzustand kann durchaus in einem
größeren Singulettbereich liegen. Die Zentrumsregel N gleich 0 modulo 4
ist nur eine notwendige Bedingung für Singuletts, kein Auswahlprinzip
für N=64.

Der Beweis für k=2 benötigt keinen riesigen Matrixkern. Die Bosonseite
zerfällt nach der symmetrischen Cauchy-Identität in

\[\adjustbox{max width=\linewidth}{$\displaystyle 
\operatorname{Sym}^2(10\otimes6)=
(1,1)\oplus(54,1)\oplus(1,20')\oplus(54,20')\oplus(45,15).
$}\]

Die Dimensionen 1, 54, 20, 1080 und 675 ergeben zusammen 1830. Die
vierte äußere Potenz der Fermionseite wird durch die fünf Partitionen
von vier beschrieben. Für die benötigten Überschneidungen reichen drei
Spinor-Schurfunktoren:

{\def\LTcaptype{none} % do not increment counter
\begin{longtable}[]{@{}
  >{\raggedright\arraybackslash}p{(\linewidth - 4\tabcolsep) * \real{0.3000}}
  >{\raggedright\arraybackslash}p{(\linewidth - 4\tabcolsep) * \real{0.3000}}
  >{\raggedleft\arraybackslash}p{(\linewidth - 4\tabcolsep) * \real{0.4000}}@{}}
\toprule\noalign{}
\begin{minipage}[b]{\linewidth}\raggedright
Bosontyp
\end{minipage} & \begin{minipage}[b]{\linewidth}\raggedright
Spinor-Schurfunktor auf der Fermionseite
\end{minipage} & \begin{minipage}[b]{\linewidth}\raggedleft
Anzahl gemeinsamer Typen
\end{minipage} \\
\midrule\noalign{}
\endhead
\bottomrule\noalign{}
\endlastfoot
(1,1) & Sym hoch 4 von 16 & 0 \\
(54,1) & Sym hoch 4 von 16 & 1 \\
(1,20') & S mit Partition (2,2) von 16 & 1 \\
(54,20') & S mit Partition (2,2) von 16 & 1 \\
(45,15) & S mit Partition (3,1) von 16 & 1 \\
\end{longtable}
}

Die Charaktere wurden aus den fünf Konjugationsklassen von S4 mit
ganzzahliger Arithmetik aufgebaut, durch eine unabhängige
Jacobi-Trudi-Rechnung kontrolliert und mit der Weyl-Alternierung
ausgewertet. Die volle äußere Potenz hat 635376 Dimensionen; der
gemeinsame Singulettbereich hier nur vier. Die Berechnung des anderen
Workers über den vollständigen Gewicht-Null-Raum kommt unabhängig auf
denselben Wert.

\subsubsection{3.2 Nicht nur gezählt: die vier Zustände und ihre
Kopplungen}\label{nicht-nur-gezuxe4hlt-die-vier-zustuxe4nde-und-ihre-kopplungen}

In der durch eine Bosonparität äquivalenten Lochdarstellung sei
\(v_k=(T_+)^kF\). Alle 108240 von null verschiedenen Koeffizienten von
\(v_2\) wurden neu berechnet, einschließlich der Boson-Fakultätsnormen.
Das ergibt \(\|v_2\|^2=439680\).

Sei \(P_R\) der orthogonale Projektor auf den jeweiligen inneren Typ der
Zwei-Boson-Seite. Dann sind die vier Vektoren \(P_Rv_2\) nicht null,
orthogonal und bilden wegen der bewiesenen Multiplizität eins die
gesamte Singulettbasis auf Stufe zwei. Ihre Normen und die Kopplungen
vom normierten Ein-Boson-Zustand lauten:

{\def\LTcaptype{none} % do not increment counter
\begin{longtable}[]{@{}
  >{\raggedright\arraybackslash}p{(\linewidth - 4\tabcolsep) * \real{0.2727}}
  >{\raggedleft\arraybackslash}p{(\linewidth - 4\tabcolsep) * \real{0.3636}}
  >{\raggedleft\arraybackslash}p{(\linewidth - 4\tabcolsep) * \real{0.3636}}@{}}
\toprule\noalign{}
\begin{minipage}[b]{\linewidth}\raggedright
Typ R
\end{minipage} & \begin{minipage}[b]{\linewidth}\raggedleft
Normquadrat von P\_R v2
\end{minipage} & \begin{minipage}[b]{\linewidth}\raggedleft
Quadrierte Kopplung geteilt durch g hoch 2
\end{minipage} \\
\midrule\noalign{}
\endhead
\bottomrule\noalign{}
\endlastfoot
(54,1) & 17280 & 36 \\
(1,20') & 7680 & 16 \\
(54,20') & 241920 & 504 \\
(45,15) & 172800 & 360 \\
Summe & 439680 & 916 \\
\end{longtable}
}

Diese Zahlen sind abgeleitet, nicht angepasst. Die Projektoren benutzen
nur Symmetrisierung sowie die beiden vorhandenen invarianten
Spurbildungen. Die gemeinsame skalare Projektion (1,1) verschwindet
identisch.

Nach Wahl der vier Basisphasen ist der Kopplungsvektor

\[\adjustbox{max width=\linewidth}{$\displaystyle 
g\,(6,\,4,\,6\sqrt{14},\,6\sqrt{10}).
$}\]

Alle vier Richtungen sind in v2 enthalten, aber von Stufe eins wird nur
ihre eine feste Linearkombination angeregt. Drei dazu orthogonale
Kombinationen sind \textbf{nach unten} dunkel. Sie sind nicht deshalb
auch gegenüber den höheren Stufen entkoppelt.

\subsubsection{3.3 Der vollständig bestimmte
Sechszustands-Ausschnitt}\label{der-vollstuxe4ndig-bestimmte-sechszustands-ausschnitt}

Die Kompression auf alle Singuletts mit k kleiner oder gleich zwei ist
exakt äquivalent zu

\[\adjustbox{max width=\linewidth}{$\displaystyle 
P_{\leq2}HP_{\leq2}\simeq
\begin{pmatrix}
0&g\sqrt{480}&0\\
g\sqrt{480}&\Delta&g\sqrt{916}\\
0&g\sqrt{916}&2\Delta
\end{pmatrix}
\oplus 2\Delta I_3.
$}\]

Damit ist dieser Ausschnitt vollständig gelöst. \textbf{Er ist nicht
invariant:} H führt von k=2 auch nach k=3. Seine Eigenwerte sind
folglich nicht automatisch Eigenwerte der ganzen Bank. Eine räumliche
Interpretation oder zusätzliche Wechselwirkung wurde für diese Reduktion
nicht eingeführt.

Explizit besitzt die Kompression dreimal den Wert 2 Delta und die drei
reellen Nullstellen des Polynoms

\[\adjustbox{max width=\linewidth}{$\displaystyle 
E^3-3\Delta E^2+(2\Delta^2-1396g^2)E+960\Delta g^2.
$}\]

